\section{Conclusion}
\label{sec:conclusion}
% src: every statement is a summary of Sect. 3; numbers: 90 answerable questions (Sect. 2.2) and about four times the tokens (Table search, S9_gpu/RESULTS.md: 4.09 and S3_agentic/RESULTS.md: 4.19)
This paper presented a comparison of retrieval-augmented generation setups on natural Persian and English documents, in which the retrieval and selection stages vary at fixed evidence budgets while the reader and the answer prompt stay fixed. A frozen confirmatory test of the proposed pipeline is kept separate from all later analyses. A central contribution of this work is the diagnosis of where answers are lost. The frozen test on 90 answerable questions was inconclusive, because the proposed Lex-RAG pipeline missed its required accuracy margin and the test was too small to detect the differences at stake. The post hoc analyses indicate that most wrong answers arise before the reader, either because no complete evidence bundle is among the candidates or because packing drops it. A reranking or selection step helps most at small budgets, and \hybrid{} reaches the same accuracy when it is shown several times as much text. Iterative search over the full hybrid index has the best point estimate, from one run and one judge, at approximately four times the tokens. For Persian, the evidence supports normalised text, budgets set in tokens, and a dense channel when questions and documents differ in language. The rankings of close pipelines also shift when the reader is run again, when a second judge is used and when the human review changes. Each recommendation is therefore reported with a stated confidence.

Notwithstanding these findings, several limitations define the scope for future work. The test is small and was reused for all post hoc analyses, the two languages use different documents, and the benchmark was drafted and checked mainly by language-model agents. The human review was partial and made by one person, and each experiment ran once with one primary reader. Future research should therefore prioritise a larger fresh test with independent annotation, and a confirmatory comparison between a reranker and a selection step with enough questions to detect a difference of a few answers. A paired design that uses the same documents in both languages would separate the effect of the language from the effect of the sources. The iterative and graph-based methods that were only partly tested, together with additional readers, judges and rerankers, should also be evaluated under the same budget accounting. Ultimately, this work contributes a protocol, a benchmark and a set of findings, each with a stated level of confidence, that help practitioners to choose and to check a retrieval setup for Persian and English documents.
